% Generated draft: requires mathematical and source review.
\documentclass[11pt]{article}
\usepackage[margin=25mm]{geometry}
\usepackage{fontspec}
\setmainfont{Latin Modern Roman}
\usepackage{amsmath,amssymb,mathtools,unicode-math}
\setmathfont{Latin Modern Math}
\usepackage{enumitem,needspace,xurl,hyperref}
\setlength{\parindent}{0pt}
\setlength{\parskip}{5pt}
\setlength{\emergencystretch}{3em}
\newcommand{\N}{\mathbb{N}}
\newcommand{\Z}{\mathbb{Z}}
\newcommand{\Q}{\mathbb{Q}}
\newcommand{\R}{\mathbb{R}}
\newcommand{\C}{\mathbb{C}}
\newcommand{\F}{\mathbb{F}}
\newcommand{\A}{\mathbb{A}}
\newcommand{\PP}{\mathbb{P}}
\newcommand{\E}{\mathbb{E}}
\newcommand{\eps}{\varepsilon}
\newcommand{\dd}{\,\mathrm{d}}
\begin{document}
\section*{UnsolvedMath ID: 30004392}
\textbf{Source problem code:} catalogue entry
\subsection{Definitions and mathematical statement}
Let \((M,d)\) be a metric space. Its Assouad--Nagata dimension is the least integer \(n\geq 0\), if one exists, for which there is a constant \(c>0\) such that, for every \(r>0\), there is a cover \(\mathcal U_r\) of \(M\) satisfying
\[
\operatorname{diam}(U)\leq cr\quad(U\in\mathcal U_r),
\]
and every subset of \(M\) of diameter at most \(r\) meets at most \(n+1\) members of \(\mathcal U_r\). If no such integer exists, write \(\dim_{\mathrm{AN}}M=\infty\). The space is uniformly discrete if there exists \(\delta>0\) such that \(d(x,y)\geq\delta\) for all distinct \(x,y\in M\).

The supplied record asks the following two questions, with the second introduced there as a consequence of the first.

\begin{enumerate}
\item For every metric space \(M\), does
\[
\dim_{\mathrm{AN}}M=\infty
\quad\Longrightarrow\quad
\text{the travelling-salesman girth of }M\text{ is infinite}
\]
hold?
\item For every uniformly discrete metric space \(M\), does at least one of the following alternatives hold?
\[
\begin{aligned}
&\text{For every order }T\text{ on }M,\quad
  \forall k\in\Z_{\geq1},\quad
  \operatorname{OR}_{M,T}(k)=k-1;\\
&\text{There exist an order }T\text{ on }M\text{ and a constant }C>0
  \text{ such that}\quad
  \forall k\in\Z_{\geq1},\quad
  \operatorname{OR}_{M,T}(k)\leq C\ln k.
\end{aligned}
\]
Here \(C\) may depend on \(M\) and the chosen order \(T\), but not on \(k\). The original statement uses \(\ln\); the cleaned statement uses \(\log\).
\end{enumerate}

\textbf{Unresolved source qualifications.} The supplied record does not define travelling-salesman girth or \(\operatorname{OR}_{M,T}(k)\), specify the class of orders, or give the cardinality and path-length conventions underlying the ordering ratio. These expressions therefore retain their source-specific meanings as unresolved notation. In particular, the displayed range includes \(k=1\): the first alternative requires \(\operatorname{OR}_{M,T}(1)=0\), and the second requires \(\operatorname{OR}_{M,T}(1)\leq0\). A convention making these assertions meaningful is not supplied. An ordinary quotient of a prescribed visiting-path length by a shortest visiting-path length is at least \(1\) when its denominator is positive, while singleton sets may give a zero denominator; consequently, no such quotient definition or correction of the index is silently adopted here. The source also gives no exceptions or conventions for empty, singleton, or other finite spaces. Thus these are faithful source questions whose complete precise formulation requires the missing definitions and boundary conventions. The claimed consequential relationship between the questions is retained as source wording, not established here.
\subsection{Short English statement}
The source asks whether infinite Assouad--Nagata dimension forces infinite travelling-salesman girth, and whether every uniformly discrete metric space satisfies the stated alternative between maximal ordering ratios for every order and a logarithmic bound for some order. Definitions of girth and ordering ratio, including the convention at \(k=1\), are missing, so this statement-only draft requires clarification and human mathematical review.
\subsection{Sources}
\begin{itemize}
\item Supplied local UnsolvedMath registry record: permanent ID \(30004392\), title \emph{Travelling-Salesman Girth and Metric Ordering Ratios}, source problem code \texttt{catalogue entry}, in the \emph{Oberwolfach Reports Open Problems} collection. No individual proposer is identified; the proposed year is recorded as \(2020\).
\item Supplied report citation: \emph{Manifolds and Groups (2021)}, Oberwolfach Reports, DOI \texttt{10.4171/owr/2020/7}, pp.~417--463. \url{https://doi.org/10.4171/owr/2020/7}. The background describes these reports as documenting mathematical workshops and their open-problem discussions. The supplied citation uses the year \(2021\), whereas the DOI contains \(2020\); these details are retained without reconciliation.
\item Supplied literature source labeled ``Erschler--Mitrofanov 2023'': \url{https://ems.press/content/serial-article-files/43590}. No fuller bibliographic identification is supplied.
\item Supplied literature assessment, checked on \texttt{2026-08-22}: the record labels the problem ``partially solved'' and reports that infinite travelling-salesman girth is proved for important classes such as spaces weakly containing arbitrarily large cubes or spheres, while the implication from infinite Assouad--Nagata dimension remains open. The listed sources checked are the Oberwolfach report and Erschler--Mitrofanov 2023, linked above. This is supplied dated curation derived from the OWR material, not an independently verified assessment; its time-sensitive status and cited claims require independent verification.
\item The registry labels the statement \texttt{corrected\_verified} and supplies the verification note: the official report was checked, its two consecutive metric questions were retained together, and the ordering-ratio quantifiers were made explicit. This records the supplied verification claim; no source visit or independent verification is asserted in this draft.
\end{itemize}
\end{document}
